\uuid{dBi2}
\exo7id{542}
\titre{exo7 542}
\auteur{cousquer}
\organisation{exo7}
\datecreate{2003-10-01}
\isIndication{false}
\isCorrection{false}
\chapitre{Suite}
\sousChapitre{Suite définie par une relation de récurrence}
\cpge{mpsi-mp2i.nombres-reels-suites.i, pcsi.nombres-reels-suites.g, ptsi.nombres-reels-suites.g, mpsi-mp2i.nombres-reels-suites.d, pcsi.nombres-reels-suites.c, ptsi.nombres-reels-suites.c, mpsi-mp2i.derivabilite.c, pcsi.derivabilite.c, ptsi.derivabilite.c, bcpst1.suites-reelles, bcpst1.limites-continuite.c, bcpst1.derivation.b}
\module{Analyse}
\niveau{L1}
\difficulte{}

\contenu{
\question{
Étudier la suite $(u_n)$ définie par $u_0=0$ et 
$u_{n+1}=\frac{(u_n-3)^2}{4}$.
}
}
